% ---- I. Introduction. Source: JAMA introduction (paragraphs 1, 2 and 4), plus the automation
% framing and the contributions list that T-ASE expects.
\section{Introduction}

Ambulance routing for suspected stroke is a time-critical decision made under uncertainty, and protocols vary considerably across the United States. In parts of California the absence of a coordinated routing framework means patients are often transported to the nearest hospital by default, which can delay access to definitive care. Rhode Island offers a contrasting case study in severity-based routing, directing suspected large-vessel-occlusion patients to a single comprehensive stroke center, which has been shown to shorten times to treatment~\cite{jayaraman2023}. That approach is not practical in larger or denser regions, where travel distances grow and a single center would be overwhelmed. Between these two fixed rules lies a decision that can be computed: given where the patient is, how long it has been since onset, and what each reachable hospital can do, which destination maximizes the expected outcome?

The urgency is well quantified: approximately 1.9 million neurons are lost for every minute that treatment is delayed~\cite{saver2006}. Delivering the right treatment is complicated by the range of subtypes a patient may have. Large-vessel occlusion (LVO) is treated with endovascular thrombectomy (EVT), available only at a subset of hospitals, and with intravenous thrombolysis; non-LVO ischemic stroke (nLVO) is treated with thrombolysis alone; intracerebral hemorrhage requires neither, and a share of positive prehospital screens are stroke mimics. We therefore classify hospitals by what they can do on arrival: EVT-capable centers, thrombolysis-capable centers, and acute-care hospitals with no stroke certification. The subtype is unknown at pickup and becomes known only after imaging at the first hospital, so the destination choice is a decision under uncertainty whose consequences include a possible second journey.

Framed this way, prehospital stroke routing is a sequential decision problem with a well-defined state, a small discrete action set and a quantitative objective, rather than a matter of clinical judgment alone. Every quantity the decision depends on is available to a dispatch system in real time: the ambulance's position, the onset time reported by the caller, road travel times from a routing engine, and a table of hospital capabilities. The barrier to deployment is not data but a policy that can be computed inside the dispatch window and explained to the crew that applies it. This paper provides both: a Markov decision process (MDP) policy solved online, and its distillation into a decision tree that needs no computation at all.

Our contributions are as follows. (i) An MDP formulation of prehospital stroke routing that marginalizes over subtype uncertainty at the first decision, models one inter-facility transfer with published in-hospital intervals (door-to-needle, door-to-puncture and door-in-door-out), and uses road-network travel times from an open routing engine. (ii) An interpretable decision-tree policy, selected by a depth sweep on held-out patients, that recovers 99.8\% of the MDP's gain over nearest-hospital routing with a single split. (iii) An evaluation in two regions with ten replicates each, including comparison with the bypass protocols actually in use, stratification by neighborhood income and urbanicity, sensitivity to travel-time error, in-hospital intervals and the definition of EVT capability, and measured per-decision latency. The remainder of the paper reviews related work (Section~\ref{sec:related}), states the problem (Section~\ref{sec:problem}), describes the solution method (Section~\ref{sec:method}) and the experimental setup (Section~\ref{sec:setup}), reports results (Section~\ref{sec:results}) and discusses limitations (Section~\ref{sec:discussion}).
